\chapter{Détails de calibration}
\label{ann:calibration}

\section{Structure du code}
\label{ann:code_structure}

\todoF{Ajouter un arbre détaillé du dépôt avec brève description par fichier.
Source d'inspiration : \texttt{docs/PROJECT\_STATUS.md} section "Structure du
repo".}

\section{Repère base : mesure expérimentale}
\label{ann:repere_mesure}

L'origine du repère base du SO-101 (centre du pivot du \texttt{shoulder\_pan})
est mesurée à \textbf{+32 mm au-dessus de la surface de la table} sur le
poste expérimental (mesure au pied à coulisse, 2026-05-16). En conséquence,
un objet posé sur la table a son \textbf{Z} (centre) à $-32 + h/2$ mm avec
$h$ la hauteur de l'objet en mm.

Cette mesure doit être refaite si le robot est déplacé ou si la base est
changée.

\section{Décision D1 : recalage du \texttt{wrist\_roll}}
\label{ann:d1}

\todoF{Reprendre le contenu de PROJECT\_STATUS.md section D1 (problème de
la couture 0/4095 de l'encodeur 12 bits sur le wrist\_roll, solution avec
mesure du vrai centre et recalcul du Homing\_Offset).}

\section{Décision D2 : damier asymétrique 9$\times$6}
\label{ann:d2}

\todoF{Reprendre PROJECT\_STATUS.md section D2 (4-fold ambiguity du damier
7$\times$7 carré, solution avec damier asymétrique 9$\times$6).}

\section{Décision D3 : solveur hand-eye robuste}
\label{ann:d3}

\todoF{Reprendre PROJECT\_STATUS.md section D3 (pipeline solve initial +
cluster majoritaire + réalignement + rejet outliers).}

\section{Procédure de capture des données extrinsèques}
\label{ann:capture_extr}

\todoF{Détailler la procédure suivie : nombre de poses (60+ par caméra),
critère de diversité angulaire ($\geq 30$\textdegree{} entre poses), images annotées
sauvegardées dans \texttt{outputs/calibration\_images/}.}
